\section{\texorpdfstring{The positive-\(Q\) trace for the fifteen
equal-depth one-active
pairs}{The positive-Q trace for the fifteen equal-depth one-active pairs}}

\subsection{1. Scope and claim boundary}

This note treats the coefficient-sensitive part of the fifteen support
pairs selected by \texttt{\seqsplit{src/\texttt{\seqsplit{one\_active\_remaining\_structure.py}}}}. It
proves the physical sign on the seventy-five active-\(C\) equality rows
and composes the eight companion rows in Section 7. Two independent
audits passed the stopped trace and pair-level theorem. The exact
selector is frozen in
\texttt{\seqsplit{src/\texttt{\seqsplit{critical\_one\_active\_q\_trace\_certificate.py}}}}.

After relabelling the active species to \(C\), every equality row has

\[
 L_0=\{0,A+C,B+C\},\qquad
 L_1=T\subseteq\{A,B,2A,A+B,2B\}.                    \tag{1.1}
\]

The directed graph on each displayed support is arbitrary and strongly
connected, and all labelled channel rates are positive. Each of the
fifteen possible supports \(T\) contains at least one unary and at least
one quadratic complex. Propensities use falling factorials.

The finite regression reports a positive-debt word and a first-surplus
word of the same slow-before-fast depth two. That equality does
\textbf{not} give a critical reflected random walk. One
unary-to-quadratic reaction already lowers positive debt by one at depth
one. A second such reaction is needed only because the regression
insists on crossing \(D=0\) and recording a surplus exit.

\subsection{2. The exact level}

Put

\[
 M=A+B,\qquad Q=C-A-B=C-M.                            \tag{2.1}
\]

Every reaction in \(L_0\) preserves \(Q\). A reaction \(y\to z\) in
\(T\) changes it by

\[
 \Delta Q=|y|-|z|\in\{-1,0,1\}.                     \tag{2.2}
\]

Thus unary-to-quadratic reactions move \(Q\) toward the origin on a
large positive shell, while quadratic-to-unary reactions move it away.

For \(N\geq0\), the \(L_0\)-shell \(Q=N\) is

\[
 \Gamma_N=\{(a,b,N+a+b):a,b\in\mathbb N_0\}.         \tag{2.3}
\]

Its unique base is \(z_N=(0,0,N)\). The base is regenerative for the
\(L_0\)-chain: wait for a reaction sourced at \(0\), then follow the
full busy excursion until \(M=0\) again.

\subsection{3. Exact shell law and the sign before
perturbation}

The linkage \(L_0\) is weakly reversible and deficiency zero. Let
\(K_0,K_{AC},K_{BC}>0\) be its directed matrix-tree weights and set

\[
 u=K_{AC}/K_0,\qquad v=K_{BC}/K_0,\qquad s=u+v.
\]

The conditional product-form law on \(\Gamma_N\) is

\[
 \pi_N(a,b,N+a+b)
 =\frac1{Z_N}\frac{u^a v^b}{a!b!(N+a+b)!},\qquad
 Z_N=\sum_{m\ge0}\frac{s^m}{m!(N+m)!}.               \tag{3.1}
\]

For fixed \(r,t\geq0\),

\[
 \mathbb E_{\pi_N}[(A)_r(B)_t]
 =u^rv^t\frac{Z_{N+r+t}}{Z_N}.                       \tag{3.2}
\]

In particular,

\[
 \mathbb E_{\pi_N}A=\frac{u}{N}+O(N^{-2}),\qquad
 \mathbb E_{\pi_N}B=\frac{v}{N}+O(N^{-2}),           \tag{3.3}
\]

whereas every quadratic falling-factorial moment is \(O(N^{-2})\).

Let \({\cal E}_{12}\) be the set of directed channels in \(T\) from a
unary to a quadratic complex. Strong connectivity and the presence of
both degrees imply \({\cal E}_{12}\ne\varnothing\). If a unary source is
\(A\), respectively \(B\), write \(r_y=u\), respectively \(v\), and put

\[
 a_-:=\sum_{y\to z\in{\cal E}_{12}}\kappa_{yz}r_y>0. \tag{3.4}
\]

The stationary \(T\)-hazard therefore satisfies

\[
 \mathbb E_{\pi_N}{\cal L}_TQ
 =-\frac{a_-}{N}+O(N^{-2}).                           \tag{3.5}
\]

Every rate and orientation changes \(a_-\), but cannot change its strict
positivity. Quadratic-to-unary channels contribute only to the
\(O(N^{-2})\) term. Hence no choice of positive rational rates can
reverse the leading sign.

\subsection{4. Regenerative perturbation with all reactions
retained}

Let \(\Lambda>0\) be the aggregate rate of the channels sourced at \(0\)
in \(L_0\). For the \(L_0\)-chain, a cycle from \(z_N\) consists of an
\({\rm Exp}(\Lambda)\) base wait and a busy part. Stripping the common
\(C\) molecule turns the busy part into a two-state transient
unimolecular particle system. Its per-particle clocks are multiplied by
\(C=N+M\geq N\). Strong connectivity of \(L_0\) makes the
single-particle killing chain phase type. Consequently, for every fixed
\(p\),

\[
 \mathbb E(\text{busy time})^p=O(N^{-p}),\quad
 \mathbb E\int (M)_1\,dt=O(N^{-1}),\quad
 \mathbb E\int (M)_2\,dt=O(N^{-2}).                  \tag{4.1}
\]

The last estimate includes overlapping \(0\)-source immigrants: a second
immigrant must arrive during an \(O(N^{-1})\) busy interval. More
generally, the probability of reaching \(M\geq k\) before returning to
zero has a geometric-factorial bound uniform in \(N\). This supplies
uniform integrability of every fixed endpoint and occupation moment.

Now retain every reaction in \(T\). Fix a sufficiently small,
rate-dependent \(\rho\in(0,1/4)\), and define \(\tau_N\) to be the first
of

\[
 \{\hbox{the positive return to }M=0\},\qquad
 \{M\geq\rho N\},\qquad
 \{Q\leq N/2\},\qquad
 \{Q\geq2N\}.                                        \tag{4.2}
\]

The last three endpoints are retained physical states. They are genuine
multi-scale exits, not finite-box truncations: an inactive coordinate is
of active order, the positive level has already made a macroscopic
descent, or the level has made a macroscopic upward excursion which is
charged explicitly. Equivalently, before any boundary one may stop first
at a degree-changing \(T\)-reaction and append the physical return or
boundary hit; no clock is suppressed during the append.

Only unary-to-quadratic reactions can increase \(M\). Their total rate
is at most linear in \(M\); quadratic reactions cannot increase \(M\).
Here is a uniform comparison which also retains the degree-preserving
quadratic clocks. Strip the common \(C\) from \(AC,BC\), and let \(R\)
be the transient two-state single-particle generator, with absorption at
\(0\). Since \(-R\) is a nonsingular M-matrix, there are \(r_A,r_B>1\)
and \(c>0\) such that the multiplicative particle corrector

\[
 G(a,b)=r_A^a r_B^b
\]

has top drift at most \(-cCMG\). Every quadratic degree-preserving
reaction changes \(G\) by a bounded factor and contributes at most
\(K M^2G\); every extra \(0\)-immigration or unary-to-quadratic reaction
contributes at most \(K(1+M)G\). Choose the rate-dependent \(\rho>0\) in
(4.2) so small that \(KM^2\le cNM/4\) for \(M\le\rho N\). The stopped
multiplicative drift is then strictly negative for \(1\le M<\rho N\),
after increasing the finite threshold in \(N\).

One explicit choice is \(w=(-R)^{-1}{\bf1}>0\) and
\(r_i=1+\varepsilon w_i\). For small \(\varepsilon>0\), the logarithmic
drift bracket of one phase-\(i\) particle is
\(\varepsilon(Rw)_i+O(\varepsilon^2)\le-\varepsilon/2\), which gives the
displayed product drift after summing over particles.

The corresponding linear corrector \(W=w_AA+w_BB\) gives the block-time
bound directly. Suppress only the two slow creation clocks until the
next particle absorption. Then \(M\) is constant, the quadratic
label-change error is at most \(KM^2\), and \({\cal L}W\le-cNM\) below
the chosen \(\rho N\). Since \(W\asymp M\), stopped Dynkin yields
\({\mathbb E}(\text{next absorption time})\le C/N\), uniformly in the
current label composition.

Equivalently, expose successive phase-type particle-absorption blocks.
Even after the quadratic label changes, a block has mean duration at
most \(C/N\). Its aggregate extra-immigration/unary-birth rate is at
most \(K(1+M)\). With \(\rho\) as above, the probability of one more
creation before the next absorption is bounded by a fixed \(q<1/4\),
uniformly below the boundary; shrink the rate-dependent \(\rho\) once
more to ensure this. At block endpoints, \(M\) is therefore dominated by
the nearest-neighbor walk with up probability \(q\) and down probability
\(1-q\), killed at zero. Equivalently its family tree is dominated by a
subcritical Galton--Watson total progeny. The Catalan/path multiplicity
is included in this comparison, which gives

\[
 \mathbb P\{\text{at least }j\text{ creation blocks before return}\}
 \le Ce^{-cj}.
\]

Reaching \(M=\rho N\), or changing \(Q\) by a fixed fraction of \(N\)
without first returning to \(M=0\), requires order \(N\) creation
blocks. The resulting bound is \(Ce^{-cN}=O(N^{-k})\) for every fixed
\(k\). This supplies the claimed geometric-factorial comparison rather
than inferring it merely from a pointwise negative drift.

The implication for the two \(Q\)-boundaries is exact. Before return to
zero, let \(I_0\) count extra \(0\)-immigrations, \(U_T\) unary-to-
quadratic events, \(D_T\) quadratic-to-unary events, and \(D_0\) top
particle absorptions. Then

\[
 M=1+I_0+U_T-D_T-D_0,qquad Q-N=-U_T+D_T.
\]

Thus \(Q\le N/2\) forces \(U_T\ge N/2\), while \(Q\ge2N\), together with
\(M\ge0\), forces \(I_0\ge N-1\). Either boundary therefore requires
order \(N\) of the slow creation blocks controlled above.

Together with the phase-type \(L_0\) killing blocks, this gives, for all
sufficiently large \(N\),

\[
 \sup_N\mathbb E_{z_N}\tau_N^p<\infty,\qquad
 \sup_N\mathbb E_{z_N}
 \left(1+\frac{M_{\tau_N}+C_{\tau_N}}N\right)^p<\infty
                                                               \tag{4.3}
\]

for every fixed \(p\). A direct construction is to expose consecutive
single-particle killing blocks. During one block a unary birth has
probability \(O(N^{-1})\); after such a birth the same estimate restarts
from at most two particles. Quadratic and degree-preserving reactions
may change the particle labels, but do not increase \(M\), and every
label has a finite top path to \(0\). Iteration up to \(M=\rho N\) gives
a subcritical geometric bound. Reaching that boundary requires order
\(N\) successful immigration or unary-birth races against order-\(N\)
top clocks. Hence, for every \(k<\infty\),

\[
 \mathbb P_{z_N}\{M_{\tau_N}\geq\rho N
                 \text{ or }Q_{\tau_N}\notin(N/2,2N)\}
 =O(N^{-k}).                                         \tag{4.4}
\]

The boundary overshoot is bounded because every population jump is
bounded, and (4.3) makes every fixed polynomial endpoint cost on (4.4)
negligible. This argument retains the physical quadratic clocks; they
only insert label changes or deaths and never increase \(M\).

The probability of two or more \(T\)-reactions in the same busy
excursion is \(O(N^{-2})\). The probability that the first
degree-changing \(T\)-reaction is quadratic-to-unary is also
\(O(N^{-2})\), by (4.1). Kac's cycle identity applied to (3.3) gives the
first-order coefficient of the unary-to-quadratic event. Since
\(\pi_N(z_N)=1+O(N^{-1})\),

\[
 \begin{aligned}
 \mathbb E_{z_N}(Q_{\tau_N}-N)
   &=-\frac{a_-}{\Lambda N}+O(N^{-2}),\\
 \mathbb E_{z_N}(Q_{\tau_N}-N)^2
   &= \frac{a_-}{\Lambda N}+O(N^{-2}).
\end{aligned}\tag{4.5}
\]

For completeness, if \(\sigma_N\) denotes the unperturbed positive
\(L_0\)-return cycle, continuous-time Kac gives

\[
 {\mathbb E}_{z_N}\int_0^{\sigma_N} f(X_t)\,dt
 =\frac{{\mathbb E}_{\pi_N}f}
        {\pi_N(z_N)\Lambda}.                         \tag{4.5a}
\]

Applied to the aggregate unary-to-quadratic propensity, (3.2) makes the
right side \(a_-/(\Lambda N)+O(N^{-2})\). Total unary \(T\)-occupation
is \(O(N^{-1})\), total quadratic occupation is \(O(N^{-2})\), and after
one unary event the remaining busy occupation is again \(O(N^{-1})\).
Compensation therefore makes two \(T\)-events an \(O(N^{-2})\) event.
Coupling the full and unperturbed cycles until their first \(T\)-event
changes the Kac coefficient only by \(O(N^{-2})\), which proves both
lines of (4.5).

These are full-reaction stopped-cycle expansions; (4.4) is included in
their remainders. A unary conversion followed by a degree-changing
reaction, a reversal after a unary birth, and a boundary carrier
launched near the return time all contain a second lower-during- fast
contest and are part of the \(O(N^{-2})\) remainder.

\subsection{5. Strict quadratic trace
drift}

At an ordinary base endpoint \(M=0\), one has \(Q=C\geq0\); the boundary
endpoints in (4.4) have negligible quadratic contribution. Combining the
two lines of (4.5),

\[
 \begin{aligned}
 \mathbb E_{z_N}\bigl[Q_{\tau_N}^2-N^2\bigr]
 &=2N\,\mathbb E_{z_N}(Q_{\tau_N}-N)
   +\mathbb E_{z_N}(Q_{\tau_N}-N)^2\\
 &=-\frac{2a_-}{\Lambda}+O(N^{-1}).
 \end{aligned}\tag{5.1}
\]

Therefore there are \(N_0<\infty\) and \(\delta>0\), depending on the
oriented network and rates, such that

\[
 \mathbb E_{z_N}\bigl[Q_{\tau_N}^2-N^2+\delta\tau_N\bigr]
 \leq-\delta,\qquad N\geq N_0.                       \tag{5.2}
\]

This is the required positive-\(Q\) physical-time sign. In particular,
the equality of the two canonical surplus-word exponents cannot produce
a null or transient reflected level on these rows.

\subsection{6. Relation to reflected
debt}

On a fast \(L_0\) carrier, entries add one to both \(D\) and \(M\), top
exits at positive debt subtract one from both, and top conversions
preserve both. Thus \(D-M\) is unchanged by the carrier dynamics until a
surplus exit. A unary-to-quadratic \(T\)-reaction raises \(M\) by one
and hence lowers positive base debt by one after the carrier is drained.
It occurs at one slow-before-fast contest. A quadratic-to-unary reaction
can create one unit of base debt, but first requires two simultaneous
cofactors and is one order rarer. The canonical depth-two surplus word
from \(D=1\) is just two consecutive depth-one debt reductions.

\subsection{7. A single proper potential for all eighty-three
rows}

The quadratic trace need not be glued to the rest of the atlas as a
separate potential. Fix any network-dependent linear corrector
\(\ell\in\mathbb R^3\), add a constant \(K\), and put

\[
 F_\ell(x)=K+\sum_i\log(x_i!)+\ell\cdot x,\qquad
 V(x)=F_\ell(x)^2,                                   \tag{7.1}
\]

where \(K\) is large enough that \(F_\ell\geq1\). This \(V\) is proper.
On the base \(z_N\),

\[
 F_\ell(z_N)=N\log N+O(N),\quad
 F_\ell(z_{N-1})-F_\ell(z_N)=-\log N+O(1).           \tag{7.2}
\]

The event \(Q_{\tau_N}=N-1\) has probability
\(a_-/(\Lambda N)+O(N^{-2})\); every nonzero event other than this
leading one has total probability \(O(N^{-2})\), with the uniform
endpoint moments from (4.3)--(4.4). Therefore

\[
 \mathbb E_{z_N}\{V(X_{\tau_N})-V(z_N)\}
 =-\frac{2a_-}{\Lambda}(\log N)^2+O(\log N).         \tag{7.3}
\]

In particular, the same proper \(V\), rather than \(Q^2\), pays the
physical duration with a diverging margin. Beginning at any fixed
inactive population on the same positive-\(Q\) shell, first run the full
physical cleanup to \(M=0\). The no-\(T\) cleanup lowers \(F_\ell\) by
\(M\log N+O_M(1)\); a \(T\)-interruption has probability \(O_M(N^{-1})\)
and has all fixed moments. Appending the base cycle preserves (7.3),
with constants allowed to depend on that fixed inactive state.

The dependence on the inactive state is not a hidden finite-box
assumption. In a bad-sequence proof, a genuinely one-active sequence has
eventually fixed inactive coordinates. If either inactive coordinate is
unbounded, the exact source-rate subsequence has at least two active
coordinates and is routed to the corresponding multi-active descriptor.

The fifteen pairs have eighty-three feasible failures. The seventy-five
rows treated above are the five capped active-\(C\) rows on each pair.
The remaining eight rows consist of

\begin{itemize}
\tightlist
\item
  six direct-enabled rows with no positive-debt base; and
\item
  two zero-source seed rows with zero-contest service and strictly
  deeper debt creation.
\end{itemize}

After normalizing the active species to \(C\) and exchanging inactive
\(A/B\) where needed, the executable payload gives exactly the following
four templates; every row has cap \((A,B)=(0,0)\):

\[
\begin{array}{c|c|c}
L^{(1)}&L^{(2)}&\text{multiplicity/access}\\ \hline
\{0,AB,AC\}&\{B,2B,BC\}&2\ \text{zero-source}\\
\{0,AB,AC\}&\{C,2B,BC\}&2\ \text{direct}\\
\{0,AB,AC\}&\{B,C,2B,BC\}&2\ \text{direct}\\
\{0,AB,AC\}&\{B,C,BC\}&2\ \text{direct}.
\end{array}                                           \tag{7.4a}
\]

For a direct row, mark an enabled active-degree-one source and follow a
shortest top path to its first lower target. Strong connectivity makes
the fast path phase type. A completed path has

\[
 \Delta F_\ell=-\log X+O(1),                         \tag{7.4}
\]

whereas a lower-during-fast interruption has probability \(O(X^{-1})\)
and positive \(F_\ell\)-cost \(O(\log X)\). Endpoint and duration
moments are uniform from the fixed inactive start. Hence

\[
 \mathbb E\Delta V
 =2F_\ell\,\mathbb E\Delta F_\ell
   +\mathbb E(\Delta F_\ell)^2
 \leq-cF_\ell\log X                                  \tag{7.5}
\]

for all sufficiently large \(X\).

More explicitly, in each direct template the initially enabled source is
\(C\). A \(C\)-edge to \(B\) or \(2B\) is already a strict active exit.
A \(C\)-edge to \(BC\) launches the stripped open phase on \(\{0,B\}\):
on the \(X\)-clock, \(C\to BC\) is immigration and \(BC\to C\) is linear
death. It is countable, not a finite carrier. Kill this
immigration--death phase at the first \(C\)- or \(BC\)-source edge to
\(B\) or \(2B\). Strong connectivity makes that killing accessible from
its unique closed fast class, so the killed phase has a uniform
exponential absorption time and exponential \(B\)-endpoint moments. All
competitors from the first linkage or from \(B,2B\) have order at most
polynomial in this exponentially controlled \(B\) and act for physical
time \(O(X^{-1})\). Their aggregate interruption probability and
propensity-times-log cost are therefore \(O(X^{-1})\). The successful
endpoint loses one \(C\), so (7.4)--(7.5) follow with all clocks
retained.

For either zero-source row, start at the actual target of the \(0\)-edge
and run the finite lower/top return prefix with all clocks retained. A
failed seed attempt returns to the exact capped face; the number of
attempts is geometric. At the first enabled unpaired top exit, (7.4)
holds. Positive interrupted-carrier cost is again \(O(\log X/X)\). Thus
(7.5) holds for these two rows as well.

For clarity, the two possible seed targets are \(AB\) and \(AC\). A seed
at \(AB\) immediately enables the order-\(X\) \(BC\) carrier in the
second linkage. A seed at \(AC\) either returns to \(0\), an exact
failed reset, or reaches \(AB\) through the first strongly connected
linkage and then launches that same \(BC\) carrier. Strong connectivity
makes the number of resets geometric with a rate-dependent positive
success parameter. During an \(AC\) or \(BC\) carrier, every noncarrier
clock is order one, so the total exceptional probability is
\(O(X^{-1})\); its positive squared- factorial cost is
\(O(F_\ell\log X/X)\), negligible beside the successful
\(-cF_\ell\log X\) term. This verifies the two seed rows rather than
appealing to the withdrawn uniform old-debt lemma.

Finally, along every feasible descriptor of these pairs which is not one
of the eighty-three failed rows, the Anderson--Kim top-source argument
gives the usual strict factorial drift for \(F_\ell\). For bounded
jumps,

\[
 {\cal L}V=2F_\ell{\cal L}F_\ell+
 \sum_r\lambda_r(\Delta_rF_\ell)^2.                  \tag{7.6}
\]

On a passing source-rate sequence, the second term is
\(o(F_\ell|{\cal L}F_\ell|)\): every \(\Delta_rF_\ell=O(\log(2+|x|))\),
while \(F_\ell\asymp |x|\log(2+|x|)\), and the strict top-source
logarithmic ratio diverges. Explicitly, if \(\alpha_n\) is a maximal
source propensity and \(g_n\to\infty\) the certified logarithmic exit
gap, the ordinary tier estimate gives
\({\cal L}F_\ell\le-c\alpha_ng_n\), whereas

\[
 \sum_r\lambda_r(\Delta_rF_\ell)^2
 \le C\alpha_n\log^2(2+|x_n|).
\]

The ratio to \(F_\ell\alpha_ng_n\) is at most
\(C\log(2+|x_n|)/\{|x_n|g_n\}=o(1)\). Thus the passing-region drift also
remains strict for the same \(V\).

At a boundary in (4.2), the exact descriptor has at least two active
coordinates, has already lost a fixed fraction of \(Q\), or is a
super-polynomially rare larger positive-\(Q\) restart. In the first case
every affine-feasible descriptor of these fifteen pairs passes the
top-source test: by construction they are a subset of the exact
1,227-pair selector whose nonempty feasible failure set is entirely
one-active. In the last case the same critical episode restarts at the
new level. The squared-factorial estimate (7.6) charges the promoted
endpoint, while (4.4) makes the upward-restart \(V\)-cost summable.

These estimates give the following independently audited exact
composition statement.

\begin{quote}
\textbf{Theorem 7.1.} Give each linkage of any of the fifteen selected
support pairs an arbitrary strongly connected orientation and arbitrary
positive channel rates. Then every closed irreducible class is positive
recurrent.
\end{quote}

Indeed, if no finite exceptional set admitted a random-time
\(V+\)duration inequality, choose a bad state sequence in one closed
class and pass to its exact D-tier/source-rate descriptor. An infeasible
descriptor cannot occur in that affine class. A passing descriptor is
excluded by (7.6); one of the seventy-five equality rows is excluded by
(7.3); and one of the eight companion rows is excluded by (7.5). This
exhausts the finite descriptor table. The episode endpoints have the
moments used above, so the statewise selector gives

\[
 \mathbb E_x[V(X_\tau)-V(x)+\tau]\leq-1              \tag{7.7}
\]

outside a finite classwise set. Only source-molecularity-zero or
source-molecularity-one reactions can increase total population, at an
affine rate, so the chain is nonexplosive. Random-time Foster then gives
finite mean return to the finite set.

\subsection{8. Audit status and universal
interface}

The argument above explains what a universal 1,227-pair theorem would
have to prove. For every arbitrary strong orientation, classify
reflected creation and debt-reduction histories by their first nonzero
slow-before-fast coefficient. Strict kinetic separation plugs directly
into the common squared-factorial potential \(V\). Equality cannot be
settled by a Hamilton-cycle exponent table; for the exact fifteen-pair
equality family it is resolved by the invariant \(Q\) and the shell
trace (4.3). Any other equality family would require its own coefficient
identity.

No assertion in this note promotes any pair outside these exact fifteen.
The canonical regression is not an arbitrary-orientation theorem, a box
exit is not automatically a promotion, and no count is inferred from
kinetic depth alone.

Two independent audits checked (4.3)--(4.5a), the super-polynomial
boundary restart, all four companion templates, the squared-potential
gluing estimate (7.6), and the exact prior-branch disjointness. The
certified update is

\[
 (1835,187)\longmapsto(1820,187).
\]

The analytic and pair-level flags are true only for these fifteen pairs.
The global T3-2 flag remains false.
